% ===========================================================================
\section{Evaluation Parameters}
% Source material: docs/evaluation-parameters.md
% The rubric explicitly grades "defined evaluation parameters", so this section
% carries real weight. Every parameter needs a measurement method, not just a
% name.
% ===========================================================================

The design is evaluated against five families of parameter: cost, performance,
scalability, security, and usability. Each parameter names what is measured, how
it is measured, and what result counts as success. Security parameters instead
name an adversary that must fail.

The targets are hypotheses. Nothing has been built, measured or deployed. A
missed target will be reported with its measured value and the reason it was
missed.

\subsection{Cost}

The same bytecode and calldata consume the same gas on every machine, so A1
measures the design rather than the market. A2 converts that figure into rupees
using a gas price and a token price, and states both, so the result can be
recomputed when prices move. A3 applies three fee regimes to the same A1 figure,
for Ethereum mainnet, Polygon PoS mainnet and Polygon Amoy, and reports a band
rather than a single number. One measurement serves all three, because identical
bytecode burns identical gas.

Cost is unlikely to be the limit that binds first. Throughput is, and C3 locates
it.

\begin{table}[htbp]
  \centering
  \caption{Cost parameters.}
  \label{tab:eval-cost}
  \small
  \begin{tabular}{@{}lL{5.3cm}L{5.1cm}L{3.9cm}@{}}
    \toprule
    \textbf{ID} & \textbf{Parameter}        & \textbf{Measurement}                         & \textbf{Target}       \\
    \midrule
    A1          & Gas per operation         & Gas snapshot per contract entry point & Ballot under 400k gas \\
    A2          & Cost per vote in INR      & A1 times gas price times token price         & Under INR 0.50        \\
    A3          & Cost across fee markets   & A1 priced under three fee regimes            & A band, not a point   \\
    A4          & On-chain vs IPFS metadata & Inline bytes vs a 32-byte identifier         & Order-of-magnitude gap \\
    A5          & Proxy deployment saving   & Direct deploy vs minimal proxy clone         & Over 90\% saving      \\
    \bottomrule
  \end{tabular}
\end{table}

The A1 budget is tight, since one Groth16 pairing check costs roughly
200k gas. An overshoot means anonymity on the EVM costs more than budgeted.

\subsection{Performance}

B1 and B2 report the 95th percentile beside the median: the voter who waited
is the one who abandons the ballot. B2 names its device class, since the
project does not choose the voter's hardware. On a phone the binding limit is
memory, not speed: artifacts run through snarkjs under a mobile WASM ceiling,
so B2 records artifact size and peak memory beside the timing and counts an
outright failure as a failure. The Section~6 gate turns on B2, since
in-browser proving is the principal technical risk: a 95th percentile in tens
of seconds, or a device that cannot finish a proof, fails it. It fires one
week after M1 lands, or immediately if M1 has not landed by 18 October. The
fallback is not free: server-side proving moves the identity secret onto the
operator the design exists to remove.

\begin{table}[htbp]
  \centering
  \caption{Performance parameters.}
  \label{tab:eval-perf}
  \small
  \begin{tabular}{@{}lL{5.3cm}L{5.1cm}L{3.9cm}@{}}
    \toprule
    \textbf{ID} & \textbf{Parameter}           & \textbf{Measurement}                   & \textbf{Target}                \\
    \midrule
    B1          & End-to-end vote latency      & Wall clock, click to confirmation      & 95th percentile under 10 s     \\
    B2          & Client proof generation time & Timed in-browser, by device class      & Under 5 s on a mid-range phone \\
    B3          & On-chain verification cost   & Gas attributable to proof verification & Isolate from base cost         \\
    B4          & Throughput                   & Ramp 1 to 200 ballots through the relayer & Knee of the curve \\
    B5          & Indexer lag                  & \texttt{VoteCast} to audit dashboard   & Under 3 s                      \\
    \bottomrule
  \end{tabular}
\end{table}

\subsection{Scalability}

C1 is the central claim, stated precisely enough to be refuted. It is narrower
than ``gas does not grow with the electorate''. Two operations are constant in
the roll size: publishing a roll, and casting a tier 3 ballot. The tree is
built off-chain, and only its 32-byte root, the roll size and the leaf-set
identifier go on-chain. C1 measures gas at rolls of $10^1$ to $10^5$.

Registration gas is constant by construction: a root, a roll size and a
content identifier are written whatever the roll holds. The prediction is
exact equality at all five sizes, so any variation means an O(n) term entered
the implementation.

Tier 3 \texttt{castVote} gas varies by under 1\% from $10^1$ to $10^5$, with
no monotone trend, because a Groth16 proof is one pairing check at every depth
and the path is a private witness. A larger spread or a steady rise refutes
this.

Tier 2 rises with the logarithm of the roll, because it sends the path in
calldata: one 32-byte sibling a level at 16 gas a non-zero byte plus a hash,
about 500 to 600 gas a level. Depth runs from 4 at $10^1$ to 17 at $10^5$, a
predicted rise of 7k to 8k gas. A flat line refutes that model; a much steeper
rise means the verification loop is not doing what this design assumes.

Tier 3's flat line is expected to sit above tier 2's rising line throughout:
zero knowledge removes the logarithmic factor, not the constant, and that gap
is the price of anonymity. The registration saving belongs to the Merkle
commitment, since tier 2 publishes a root too.

\begin{table}[htbp]
  \centering
  \caption{Scalability parameters.}
  \label{tab:eval-scale}
  \small
  \begin{tabular}{@{}lL{5.3cm}L{5.1cm}L{3.9cm}@{}}
    \toprule
    \textbf{ID} & \textbf{Parameter}         & \textbf{Measurement}                                    & \textbf{Target}                          \\
    \midrule
    C1          & Gas vs voter-roll size     & Rolls from $10^1$ to $10^5$: registration, tier 2 ballot, tier 3 ballot & Registration and tier 3 flat; tier 2 logarithmic \\
    C2          & Proof time vs tree depth   & B2 across the same roll sizes                           & Logarithmic, not linear                  \\
    C3          & Concurrent-voter ceiling   & Load until failure                                      & Name the bottleneck                      \\
    C4          & Candidate-count scaling    & 2, 10, and 50 candidates                                & Constant per ballot                      \\
    C5          & Off-chain cost vs roll size & Tree build, leaf-set size, in-browser reconstruction   & Practical in a browser at $10^5$         \\
    \bottomrule
  \end{tabular}
\end{table}

C2 checks that the cost did not move to the client. C5 says where the O(n)
work went: onto each auditor, and onto the voter. The booth downloads the leaf
set and rebuilds the tree locally, which keeps the registrar from learning who
is voting and sits inside B1's clock. A booth that cannot rebuild at $10^5$
would force a registrar-supplied authentication path, which Section~4 rejects
on privacy grounds. C2 and C5 are where this case is most likely to fail.

\subsection{Security and correctness}

Each property is paired with an adversary that must fail, and each negative
test asserts the named revert reason.

\begin{table}[htbp]
  \centering
  \caption{Security and correctness parameters.}
  \label{tab:eval-sec}
  \small
  \begin{tabular}{@{}lL{5.1cm}L{9.7cm}@{}}
    \toprule
    \textbf{ID} & \textbf{Property}        & \textbf{How it is demonstrated}                                   \\
    \midrule
    D1          & Double-vote resistance   & Replaying a nullifier must revert, as must a ballot past the roll size; fuzzed over collisions \\
    D2          & Eligibility enforcement  & A valid proof against the wrong root must revert; ceremony artifacts checked by hash \\
    D3          & Ballot secrecy           & Deanonymization over all chain data; operator condition and anonymity set reported beside it \\
    D4          & Individual verifiability & Voter finds their commitment in the published leaf set, then their own ballot \\
    D5          & Universal verifiability  & Recompute the root and the tally with the project's backend switched off \\
    D6          & Authority containment    & An observer closes voting and finalizes; the authority cannot vote or alter a frozen root \\
    D7          & Static analysis          & Findings by severity, each with a written disposition             \\
    D8          & Test coverage            & Line coverage over the contract sources, target 90\%               \\
    D9          & Reorganization safety    & Simulate a chain reorganization; the indexer must self-heal       \\
    \bottomrule
  \end{tabular}
\end{table}

Groth16 needs a circuit-specific trusted setup, so D2 holds only if
Semaphore's phase-two ceremony was not compromised in its entirety. The
project uses the published artifacts and checks their hashes in the test
suite. Had every contributor colluded, forged membership proofs would verify
undetectably; the roll-size invariant bounds what they do to the count but
cannot detect one.

D3 is conditional. Against the chain nothing should be linkable, since no
on-chain value ties an address to a ballot. Against the operator, secrecy
holds only when the relayer is run separately from the registrar, or when the
voter submits directly, because the relayer reads the candidate choice and
sees the network address. D3 also reports the anonymity set size. In a roll of
a few dozen voters, with a public running tally and ordered, timestamped
ballots, it is small. A failure in D1, D2 or D6 is a design defect and
blocks the live demonstration.

\subsection{Usability and adoption}

These parameters are behavioural: each counts what a participant does or
holds. Section~3 states the risk as NFR13, since a voter who abandons the
ballot leaves turnout that reads as apathy.

\begin{table}[htbp]
  \centering
  \caption{Usability parameters.}
  \label{tab:eval-use}
  \small
  \begin{tabular}{@{}lL{5.3cm}L{5.1cm}L{3.9cm}@{}}
    \toprule
    \textbf{ID} & \textbf{Parameter}           & \textbf{Measurement}                             & \textbf{Target}        \\
    \midrule
    E1          & Steps to cast a ballot       & User actions from landing page to receipt        & At most 4              \\
    E2          & Voter gas burden             & POL a voter must hold to register and vote       & Zero                   \\
    E3          & Task success rate            & 10 users unfamiliar with cryptocurrency, unaided & Over 80\%              \\
    E4          & Time to first vote           & Wall clock on the no-wallet path                 & Under 5 minutes        \\
    E5          & Returning on a second device & A voter enrolled in one browser votes in another & A working path exists  \\
    \bottomrule
  \end{tabular}
\end{table}

E2 is zero because a voter who must first acquire POL narrows the electorate
to people who already hold cryptocurrency. A relayed ballot installs no
wallet, so E4 times wallet installation separately, as the cost of direct
submission. That is no general remedy against the relayer, because a voter
holding zero POL cannot take it. A relayer that reads the choice and declines
to forward it alters the result for exactly that population. Section~4 bounds
this with a signed receipt rather than removing it. E5 covers the run E1 to E4
miss: a cleared browser or a shared laboratory machine leaves a voter enrolled
in a frozen root who cannot prove it, and after the freeze there is no remedy.

\subsection{Measurement methodology}

The harness lives in \texttt{bench/}: raw output committed as CSV, graphs
generated by script, each point repeated at least thirty times and reported as
a median with the 95th percentile.

\begin{itemize}
  \item Gas (A1, A4, A5, C1, C4): local \texttt{forge snapshot}; gas is
        deterministic given bytecode and calldata.
  \item Throughput and ceiling (B4, C3): local Anvil, since faucet drip cannot
        fund sustained load at about 1.8 POL a ramp.
  \item Client and off-chain timings (B2, C2, C5): a developer laptop and a
        mid-range Android phone, each named.
  \item End-to-end latency (B1, B5): Polygon Amoy, sampled in Phase 4
        (9--15 November).
  \item A3: no network; it prices the A1 figure under published fee conditions.
\end{itemize}

Every number carries its controls: gas price and block number, RPC provider,
commit hash, deployed addresses, leaf-set identifier, and device for client
numbers.
